\documentclass[10pt]{beamer}
\usepackage[utf8x]{inputenc}
\usepackage[T1]{fontenc}
\usepackage[english]{babel}
\usepackage{beamerthemeCopenhagen}
\usepackage{amsmath}
\usepackage{amsfonts}
\usepackage{appendixnumberbeamer}
\usepackage{relsize}
\usepackage{graphics}
\usepackage{pgfplots}
\usepackage{tikz}
\usepackage{animate}
\usepackage{booktabs}
\usecolortheme{beaver}
\setbeamercovered{transparent}
\definecolor{formula}{rgb}{0.1, 0.1, 0.7}
\setbeamertemplate{navigation symbols}{}
\title{Quadratic Equations}

\author{}
\date{}

\begin{document}
	
	\frame{\titlepage}
\begin{frame}
	\frametitle{Lessons Overview}
	
	\setbeamercolor{block title}{bg=red!30, fg=black}
	\setbeamercolor{block body}{bg=red!10, fg=black}
	
	\begin{block}{}
		\begin{columns}
			\begin{column}{0.5\textwidth}
				\begin{itemize}
					\item \textbf{Introduction to Quadratic Equations}
					\item \textbf{Types of Equations:}
					\begin{itemize}
						\item Pure
						\item Incomplete
						\item Perfect Squares
						\item Special Trinomials
					\end{itemize}
				\end{itemize}
			\end{column}
			\begin{column}{0.5\textwidth}
				\begin{itemize}
					\item \textbf{Solving Techniques:}
					\begin{itemize}
						\item Factoring (ZPP)
						\item Completing the Square
						\item General Formula
					\end{itemize}
					\item \textbf{Discriminant Cases}
					\item \textbf{Step-by-step Exercises}
				\end{itemize}
			\end{column}
		\end{columns}
	\end{block}
\end{frame}



	
	\begin{frame}
		\frametitle{Quadratic Equations}
		
		\begin{block}{Standard Form Equation}
			A quadratic equation in standard form is written as:
			\[
			ax^2 + bx + c = 0.
			\]
			\begin{itemize}
				\item \(a\), \(b\), and \(c\) are real numbers, with \(a \neq 0\) (otherwise it wouldn't be a quadratic equation).
				\item “Standard form” simply means the terms are ordered by degree.
			\end{itemize}
		\end{block}
		
		\setbeamercolor{block body}{bg=green!20, fg=black}
		\begin{block}{Example}
			\begin{enumerate}
				\item What is the standard form of \(4 + x = -x^2\)?
				\item Is \(x^2 + 2 = 0\) in standard form?
			\end{enumerate}
		\end{block}
	\end{frame}
	
	\begin{frame}
		\frametitle{Roots of an Equation}
		
		\begin{block}{Roots of an Equation}
			To a quadratic equation like \(ax^2 + bx + c = 0\), we can associate the polynomial \(P(x) = ax^2 + bx + c\).
			
			\textit{A solution \(x_0\) is a value that, when substituted into \(P(x)\), makes the polynomial equal to zero.}
		\end{block}
		
		\setbeamercolor{block body}{bg=green!20, fg=black}
		\begin{block}{Example}
			Let’s consider the equation \(x^2 - 4 = 0\), with associated polynomial \(P(x) = x^2 - 4\).
			
			Substituting \(x = 2\):
			\[
			P(2) = 2^2 - 4 = 4 - 4 = 0
			\]
			Since \(P(2) = 0\), \(x = 2\) is a solution of the equation \(x^2 - 4 = 0\).
			
			\bigskip
			\textbf{Question:} What are the solutions to \(x^2 + 2 = 0\)?
		\end{block}
	\end{frame}
	
	\begin{frame}
		\frametitle{Square Roots and Quadratic Equation Solutions}
		
		\begin{block}{Square Root}
			The square root of a number, such as \(\sqrt{9}\), is defined as the nonnegative value whose square gives the original number. For example, \(\sqrt{9} = 3\), since \(3^2 = 9\). Yet, note that \((-3)^2 = 9\) as well.
			
			The principal square root is nonnegative; it is zero when the radicand is zero.
		\end{block}
		
		\setbeamercolor{block title}{fg=black, bg=blue!20}
		\setbeamercolor{block body}{fg=black, bg=blue!10}
		\begin{block}{Solving Quadratic Equations}
			When solving an equation like \(x^2 = 9\), we look for all numbers whose square is 9.
			
			In this case, there are two solutions: \(x = 3\) and \(x = -3\), since both \(3^2\) and \((-3)^2\) equal 9. Unlike the square root, we consider both the positive and negative roots.
		\end{block}
	\end{frame}
	
	\begin{frame}
		\frametitle{Square Roots of Negative Numbers and No Real Solutions}
		
		\begin{block}{No Real Square Root of Negative Numbers}
			The square root of a number, such as \(\sqrt{81}\), is a number that squared gives the radicand.
			
			However, there is no real square root of a negative number like \(\sqrt{-81}\), since no real number squared can yield a negative result.
		\end{block}
		
		\setbeamercolor{block title}{fg=black, bg=blue!20}
		\setbeamercolor{block body}{fg=black, bg=blue!10}
		\begin{block}{Quadratic Equations Without Solutions}
			This means that some quadratic equations have no real solution. For example, the equation \(x^2 + 2 = 0\) has no real solutions, because solving it would require computing the square root of \(-2\), which is not a real number.
		\end{block}
	\end{frame}
	
	\begin{frame}
		\frametitle{Number of Solutions of a Quadratic Equation}
		
		\begin{block}{Number of Solutions}
			A quadratic equation can have 0, 1, or 2 real solutions:
			\begin{itemize}
				\item 0 if the equation has no real roots,
				\item 1 if it has a double root,
				\item 2 if the discriminant is positive.
			\end{itemize}
		\end{block}
		
		\begin{alertblock}{Exercise}
			Determine how many solutions each equation has:
			\begin{enumerate}
				\item \(x^2 = 64\)
				\item \(x^2 = -7\)
				\item \((x + 3)^2 = 0\)
			\end{enumerate}
		\end{alertblock}
	\end{frame}
	
	\begin{frame}
		\frametitle{Quadratic Equations with One Solution}
		\fontsize{9}{11}
		
		\begin{block}{Perfect Square and Single Solution}
			A real quadratic equation has a unique real solution when its polynomial is a nonzero constant multiple of a binomial square. For a monic polynomial this takes the form:
			\[
			(x + a)^2
			\]
			The formula is:
			\[
			(x + a)^2 = x^2 + 2ax + a^2
			\]
			If a quadratic is of this form, it has one solution.
		\end{block}
		
		\setbeamercolor{block body}{fg=black, bg=green!20}
		\begin{block}{Example}
			Consider the equation \(x^2 + 4x + 4 = 0\).
			\begin{itemize}
				\item This polynomial is the square of \((x + 2)\), since:
				\[
				x^2 + 4x + 4 = (x + 2)^2
				\]
				\item Since \((x + 2)^2 = 0\), the only solution is \(x = -2\).
				\item The solution is the value that makes the binomial vanish.
			\end{itemize}
		\end{block}
	\end{frame}
	\begin{frame}
		\fontsize{8.5}{11}
		\frametitle{Pure Quadratic Equations}
		
		\begin{block}{What Are Pure Quadratic Equations?}
			Pure quadratic equations are a special type of quadratic equation where the linear term is missing. They have the form:
			\[
			ax^2 + c = 0
			\]
			In these equations, the variable appears only as a square. To solve them, isolate \(x^2=-c/a\), with \(a\ne0\). There are two real solutions if \(-c/a>0\), one if \(-c/a=0\), and none if \(-c/a<0\).
		\end{block}
		
		\setbeamercolor{block body}{fg=black, bg=green!20}
		\begin{block}{Example: Solving \(2x^2 - 8 = 0\)}
			\begin{itemize}
				\item First, isolate \(x^2\): \[2x^2 = 8\]
				\item Then, divide both sides by 2: \[x^2 = 4\]
				\item Finally, find the two numbers whose square is 4: \[x = -2 \vee x = 2 \quad (x = \pm 2)\]
			\end{itemize}
		\end{block}
	\end{frame}
	
	\begin{frame}
		\frametitle{Exercises on Quadratic Equations}
		
		\begin{alertblock}{Exercises}
			Solve the following equations:
			\begin{enumerate}
				\item \(x^2 - 5 = 0\)
				\item \(4x^2 - 1 = 0\)
				\item \(x^2 + \sqrt{2} - \sqrt{3} = 0\)
				\item \(x^2 - \sqrt{2} + \sqrt{3} = 0\)
				\item \(4x^2 - x(2 - x) + 2(x - 1) = 3\)
				\item \((x - 1)^2 + (x + 1)^2 = (x + 1)(x - 1) + 5\)
				\item Which real numbers are equal to their own square?
			\end{enumerate}
		\end{alertblock}
	\end{frame}
	
	\begin{frame}
		\frametitle{Solutions to the Exercises}
		\fontsize{9}{11}
		\setbeamercolor{block body}{fg=black, bg=green!20}
		\begin{block}{Solutions}
			\begin{enumerate}
				\item \(x^2 - 5 = 0\) \\ Solution: \(x = \pm \sqrt{5}\)
				\item \(4x^2 - 1 = 0\) \\ Solution: \(x = \pm \frac{1}{2}\)
				\item \(x^2 + \sqrt{2} - \sqrt{3} = 0\) \\ Solution: \(x = \pm \sqrt{\sqrt{3} - \sqrt{2}}\)
				\item \(x^2 - \sqrt{2} + \sqrt{3} = 0\) \\ No real solution, since \(x^2 = \sqrt{2} - \sqrt{3}\) is a square equal to a negative number.
				\item \(4x^2 - x(2 - x) + 2(x - 1) = 3\) \\ Simplifies to: \(5x^2 - 2 = 3 \Rightarrow 5x^2 = 5 \Rightarrow x = \pm 1\)
				\item \((x - 1)^2 + (x + 1)^2 = (x + 1)(x - 1) + 5\) \\ Expands to: \(2x^2 + 2 = x^2 + 4 \Rightarrow x^2 = 2\) \\ Solution: \(x = \pm \sqrt{2}\)
				\item Answer: \(0\) and \(1\) are the only real numbers equal to their own square.
			\end{enumerate}
		\end{block}
	\end{frame}
	
	\begin{frame}
		\frametitle{ZPP – Zero Product Property}
		\fontsize{9.5}{11}
		\begin{block}{Proposition}
			If the product of two numbers is zero, then at least one of the two numbers is zero:
			\[
			ab = 0 \Leftrightarrow a = 0 \vee b = 0
			\]
		\end{block}
		
		\setbeamercolor{block title}{fg=black, bg=blue!20}
		\setbeamercolor{block body}{fg=black, bg=blue!10}
		\begin{block}{Proof}
			Proof by contradiction: if denying the statement leads to a contradiction, then the statement is true.
			
			Suppose \(a \neq 0\) and \(b \neq 0\), and
			\[
			a \cdot b = 0
			\]
			
			Divide both sides by \(a\):
			\[
			\frac{1}{a} \cdot a \cdot b = \frac{1}{a} \cdot 0
			\Rightarrow b = 0
			\]
			
			But we assumed \(b \neq 0\), which is a contradiction. \(\square\)
		\end{block}
	\end{frame}
	
	\begin{frame}
		\frametitle{Incomplete Quadratic Equations}
		
		\begin{block}{What Are Incomplete Quadratic Equations?}
			Incomplete quadratic equations are equations where the constant term \(c\) is missing. They have the form:
			\[
			ax^2 + bx = 0
			\]
			To solve them, we factor out \(x\) and apply the ZPP (Zero Product Property).
		\end{block}
		
		\setbeamercolor{block body}{fg=black, bg=green!20}
		\begin{block}{Example}
			\(x^2 - 4x = 0\)
			
			Factor out \(x\):
			\[
			x(x - 4) = 0
			\]
			
			Using ZPP:
			\[
			x = 0 \quad \text{or} \quad x = 4
			\]
		\end{block}
	\end{frame}
	
	
	
	\begin{frame}
		\frametitle{Exercises – Incomplete Quadratic Equations}
		
		\begin{alertblock}{Exercises}
			Solve the following incomplete quadratic equations:
			\begin{enumerate}
				\item \(x^2 - 9x = 0\)
				\item \(2x^2 + 6x = 0\)
				\item \(x(x - 5) = 0\)
				\item \(4x^2 - 16x = 0\)
				\item \(x^2 - 2x = 0\)
			\end{enumerate}
		\end{alertblock}
	\end{frame}
	\begin{frame}
		\frametitle{Solutions to the Exercises}
		
		\setbeamercolor{block body}{bg=green!20, fg=black}
		\begin{block}{Solutions}
			\begin{enumerate}
				\item \(x^2 - 9x = 0\) \\ Solution: \(x(x - 9) = 0 \Rightarrow x = 0\) or \(x = 9\) \textcolor{formula}{[\(x = 0, 9\)]}
				
				\item \(2x^2 + 6x = 0\) \\ Solution: \(2x(x + 3) = 0 \Rightarrow x = 0\) or \(x = -3\) \textcolor{formula}{[\(x = 0, -3\)]}
				
				\item \(x(x - 5) = 0\) \\ Solution: \(x = 0\) or \(x = 5\) \textcolor{formula}{[\(x = 0, 5\)]}
				
				\item \(4x^2 - 16x = 0\) \\ Solution: \(4x(x - 4) = 0 \Rightarrow x = 0\) or \(x = 4\) \textcolor{formula}{[\(x = 0, 4\)]}
				
				\item \(x^2 - 2x = 0\) \\ Solution: \(x(x - 2) = 0 \Rightarrow x = 0\) or \(x = 2\) \textcolor{formula}{[\(x = 0, 2\)]}
			\end{enumerate}
		\end{block}
	\end{frame}
	
	\begin{frame}
		\frametitle{Quadratic Equation Exercises}
		
		\begin{alertblock}{Exercises}
			Solve the following quadratic equations:
			\begin{enumerate}
				\item \(x^2 + 2x + 1 = 0\)
				\item \(x^2 + 6x + 9 = 0\)
				\item \(x^2 + 12x + 36 = 0\)
				\item \(x^2 - 16x + 64 = 0\)
				\item \(4x^2 + 20x + 25 = 0\)
				\item \(9x^2 - 54x + 81 = 0\)
			\end{enumerate}
		\end{alertblock}
	\end{frame}
	
	\begin{frame}
		\frametitle{Solutions to the Exercises}
		
		\setbeamercolor{block body}{bg=green!20, fg=black}
		\begin{block}{Solutions}
			\begin{enumerate}
				\item \(x^2 + 2x + 1 = 0\) \textcolor{formula}{\((x + 1)^2 = 0 \Rightarrow x = -1\)}
				\item \(x^2 + 6x + 9 = 0\) \textcolor{formula}{\((x + 3)^2 = 0 \Rightarrow x = -3\)}
				\item \(x^2 + 12x + 36 = 0\) \textcolor{formula}{\((x + 6)^2 = 0 \Rightarrow x = -6\)}
				\item \(x^2 - 16x + 64 = 0\) \textcolor{formula}{\((x - 8)^2 = 0 \Rightarrow x = 8\)}
				\item \(4x^2 + 20x + 25 = 0\) \textcolor{formula}{\((2x + 5)^2 = 0 \Rightarrow x = -\frac{5}{2}\)}
				\item \(9x^2 - 54x + 81 = 0\) \textcolor{formula}{\((3x - 9)^2 = 0 \Rightarrow x = 3\)}
			\end{enumerate}
		\end{block}
	\end{frame}
	
	
	
	
	
	\begin{frame}
		\frametitle{Partial Factoring}\footnotesize
		
		\begin{block}{\footnotesize Partial Factoring}
			\footnotesize In a polynomial, we group two pairs of terms and factor each group separately:
			\[
			ax + ay + bx + by = x(a + b) + y(a + b)
			\]
			Then we factor out the common binomial:
			\[
			(a + b)(x + y)
			\]
		\end{block}
		
		\setbeamercolor{block body}{bg=green!20, fg=black}
		\begin{block}{\footnotesize Example}
			\footnotesize Let's solve the equation \(2x^2 + 4x + 3x + 6 = 0\) using partial factoring and the Zero Product Property (ZPP):
			\[
			2x(x + 2) + 3(x + 2) = 0
			\]
			Factor \((x + 2)\):
			\[
			(2x + 3)(x + 2) = 0
			\]
			Apply ZPP:
			\[
			2x + 3 = 0 \quad \text{or} \quad x + 2 = 0
			\]
			The solutions are:
			$
			x = -\frac{3}{2} \quad \text{or} \quad x = -2
			$
		\end{block}
	\end{frame}
	
	\begin{frame}
		\frametitle{Special Trinomials}
		
		\begin{block}{Special Trinomial}
			A polynomial \(P(x) = ax^2 + bx + c\) such that \(a = 1\), \(b = n_1 + n_2\), and \(c = n_1 n_2\), with \(n_1, n_2 \in \mathbb{Z}\), is called a special trinomial and can be written as:
			\[
			P(x) = (x + n_1)(x + n_2)
			\]
		\end{block}
		
		\setbeamercolor{block body}{fg=black, bg=green!20}
		\begin{block}{Example}
			An equation containing this type of polynomial can be solved using the Zero Product Property. For example, the equation \(x^2 + 7x + 12 = 0\) becomes:
			\[
			(x + 3)(x + 4) = 0
			\]
			From which the solutions are \(x = -3\) and \(x = -4\).
		\end{block}
	\end{frame}
	\begin{frame}
		\frametitle{Proof of the Special Trinomial Formula}
		
		\begin{block}{Where the Formula Comes From}
			\small Let's consider a special trinomial of the form \(P(x) = x^2 + sx + p\), where \(s = a + b\) and \(p = a \cdot b\). We can factor this trinomial using partial factoring.
			
			Start by writing the trinomial as a sum of four terms:
			\[
			P(x) = x^2 + ax + bx + ab
			\]
			Now group the terms into two pairs and factor each group:
			\[
			x(x + a) + b(x + a)
			\]
			Notice that both expressions contain the common factor \((x + a)\), which we can factor out:
			\[
			(x + a)(x + b)
			\]
			This proves that the special trinomial \(x^2 + sx + p\) can be factored as the product of \((x + a)\) and \((x + b)\).
		\end{block}
	\end{frame}
	
	\begin{frame}
		\frametitle{Special Trinomial Exercises}
		\begin{alertblock}{Exercises}
			Solve the following quadratic equations:
			\begin{enumerate}
				\item \(x^2 + 3x + 2 = 0\)
				\item \(x^2 + 5x + 6 = 0\)
				\item \(x^2 - 7x + 12 = 0\)
				\item \(x^2 - 2x - 15 = 0\)
				\item \(x^2 + 6x - 16 = 0\)
				\item \(x^2 - 5x - 24 = 0\)
			\end{enumerate}
		\end{alertblock}
	\end{frame}
	
	\begin{frame}
		\frametitle{Special Trinomial Exercises - Solutions}
		\setbeamercolor{block body}{bg=green!20, fg=black}
		\begin{block}{Solutions}
			\begin{enumerate}
				\item \(x^2 + 3x + 2 = 0\) \quad {\color{blue} \(x_1 = -1, \, x_2 = -2\)}
				\item \(x^2 + 5x + 6 = 0\) \quad {\color{blue} \(x_1 = -2, \, x_2 = -3\)}
				\item \(x^2 - 7x + 12 = 0\) \quad {\color{blue} \(x_1 = 3, \, x_2 = 4\)}
				\item \(x^2 - 2x - 15 = 0\) \quad {\color{blue} \(x_1 = 5, \, x_2 = -3\)}
				\item \(x^2 + 6x - 16 = 0\) \quad {\color{blue} \(x_1 = 2, \, x_2 = -8\)}
				\item \(x^2 - 5x - 24 = 0\) \quad {\color{blue} \(x_1 = 8, \, x_2 = -3\)}
			\end{enumerate}
		\end{block}
	\end{frame}
	
	\begin{frame}
		\frametitle{Methods for Solving Quadratic Equations}
		\setbeamercolor{block title}{bg=yellow!30, fg=black}
		\setbeamercolor{block body}{bg=yellow!10, fg=black}
		\fontsize{9.5}{11}
		\begin{block}{Summary of Methods Covered So Far}
			\begin{tabular}{|c|c|c|}
				\hline
				\textbf{Equation Type} & \textbf{Solving Method} & \textbf{Example} \\
				\hline
				Pure Quadratics & \begin{tabular}[c]{@{}c@{}}Isolate \(x^2\) and \\ extract the square root \\ to get \textbf{two} solutions.\end{tabular} & \begin{tabular}[c]{@{}c@{}}\(2x^2 = 18\) \\ \(\Rightarrow x^2 = 9\) \\ \(\Rightarrow x = \pm 3\)\end{tabular} \\
				\hline
				Spurious Quadratics & \begin{tabular}[c]{@{}c@{}}Factoring \\ and ZPP.\end{tabular} & \begin{tabular}[c]{@{}c@{}}\(x^2 - 5x = 0\) \\ \(x(x - 5) = 0\) \\ \(\Rightarrow x = 0 \, \text{or} \, x = 5\)\end{tabular} \\
				\hline
				Binomial Square & \begin{tabular}[c]{@{}c@{}}Recognize the square \\ of a binomial and solve.\\ Has \textbf{one} solution.\end{tabular} & \begin{tabular}[c]{@{}c@{}}\(x^2 + 8x + 16\)\\\((x + 4)^2 = 0\) \\ \(\Rightarrow x = -4\)\end{tabular} \\
				\hline
				Special Trinomial & \begin{tabular}[c]{@{}c@{}}Factoring \\ and ZPP.\end{tabular} & \begin{tabular}[c]{@{}c@{}}\(x^2 + x - 6 = 0\) \\ \((x + 3)(x - 2) = 0\) \\ \(\Rightarrow x = -3 \, \text{or} \, x = 2\)\end{tabular} \\
				\hline
			\end{tabular}
		\end{block}
	\end{frame}
	
	\begin{frame}
		\frametitle{Solving \((3x - 5)^2 = 16\)}
		\fontsize{9}{11}
		\begin{block}{Step-by-Step Solution}
			\begin{enumerate}
				\item Start by taking the square root of both sides:
				\[
				(3x - 5)^2 = 16 \Rightarrow 3x - 5 = \pm \sqrt{16} = \pm 4
				\]
				\item Solve both resulting equations:
				\[
				1. \ 3x - 5 = 4 \quad \text{and} \quad 2. \ 3x - 5 = -4
				\]
				\textbf{First:}
				\[
				3x = 9 \Rightarrow x = 3
				\]
				\textbf{Second:}
				\[
				3x = 1 \Rightarrow x = \frac{1}{3}
				\]
				\item Final solutions:
				\[
				x_1 = 3 \quad \text{and} \quad x_2 = \frac{1}{3}
				\]
			\end{enumerate}
		\end{block}
	\end{frame}
	
	\begin{frame}
		\frametitle{Exercises}
		\begin{alertblock}{Solve the following equations:}
			\begin{enumerate}
				\item \((3x - 5)^2 = 16\)
				\item \((x + 7)^2 = 9\)
				\item \(\left(\frac{1}{3} - \frac{1}{2}x\right)^2 = \frac{1}{4}\)
				\item \((\sqrt{2}x - 1)^2 - 32 = 0\)
				\item \(\left(-5x - \frac{2}{3}\right)^2 = \frac{1}{9}\)
			\end{enumerate}
		\end{alertblock}
	\end{frame}
	
	\begin{frame}
		\frametitle{Solutions}
		
		\setbeamercolor{block body}{bg=green!20, fg=black}
		\begin{block}{Solutions}
			\begin{enumerate}
				\item \((3x - 5)^2 = 16 \Rightarrow x = 3\) or \(x = \frac{1}{3}\)
				\item \((x + 7)^2 = 9 \Rightarrow x = -4\) or \(x = -10\)
				\item \(\left(\frac{1}{3} - \frac{1}{2}x\right)^2 = \frac{1}{4} \Rightarrow x = -\frac{1}{3}\) or \(x = \frac{5}{3}\)
				\item \((\sqrt{2}x - 1)^2 = 32 \Rightarrow \sqrt{2}x-1=\pm4\sqrt{2}\) \\ \(x=\frac{1}{\sqrt{2}}\pm4\)
				\item \(\left(-5x - \frac{2}{3}\right)^2 = \frac{1}{9} \Rightarrow x = -\frac{1}{5}\) or \(x = -\frac{1}{15}\)
			\end{enumerate}
		\end{block}
	\end{frame}
	
	\begin{frame}
		\frametitle{Completing the Square}
		
		\fontsize{9}{11}
		\begin{block}{How do we solve \(x^2 + 4x - 6 = 0\)?}
			The goal is to transform the equation \(x^2 + 4x - 6 = 0\) into a solvable form by completing the square, that is, rewriting \(x^2 + 4x\) as a perfect square.
			
			\begin{enumerate}
				\item Identify the constant needed to complete the square. For \(x^2 + 4x\), we need \(4\):
				\[
				(x + 2)^2 = x^2 + 4x + 4
				\]
				\item Add and subtract \(4\) within the equation:
				\[
				x^2 + 4x - 6 = (x^2 + 4x + 4) - 4 - 6 = (x + 2)^2 - 10
				\]
				\item Solve:
				\[
				(x + 2)^2 = 10 \Rightarrow x + 2 = \pm \sqrt{10}
				\]
				\[
				x = -2 \pm \sqrt{10}
				\]
			\end{enumerate}
			
			Final solutions: \(x = -2 + \sqrt{10}\) and \(x = -2 - \sqrt{10}\)
		\end{block}
	\end{frame}
	
	
	\begin{frame}
		\frametitle{General Completing the Square Formula}\footnotesize
		
		\begin{block}{Formula}
			In general, to complete the square for an expression of the form \( x^2 + bx \), we add and subtract \( \left( \frac{b}{2} \right)^2 \):
			\[
			x^2 + bx + \left( \frac{b}{2} \right)^2 - \left( \frac{b}{2} \right)^2 = \left( x + \frac{b}{2} \right)^2 - \left( \frac{b}{2} \right)^2
			\]
		\end{block}
		
		\setbeamercolor{block body}{bg=green!20, fg=black}
		\begin{block}{Example: Solve \( x^2 + 6x + 1 = 0 \) by completing the square}
			Identify the coefficient of \( x \), which is \( 6 \). Complete the square by adding and subtracting \( \left( \frac{6}{2} \right)^2 = 9 \):
			\[
			x^2 + 6x + 1 = (x^2 + 6x + 9) - 9 + 1 = (x + 3)^2 - 8
			\]
			So:
			\[
			(x + 3)^2 - 8 = 0 \quad \Rightarrow \quad (x + 3)^2 = 8
			\]
			\[
			x + 3 = \pm \sqrt{8} \quad \Rightarrow \quad x = -3 \pm \sqrt{8}
			\]
			Final solutions:
			\[
			x = -3 + \sqrt{8} \quad \text{or} \quad x = -3 - \sqrt{8}
			\]
		\end{block}
	\end{frame}
	
	\begin{frame}
		\frametitle{Exercises – Completing the Square}
		\begin{alertblock}{Exercises}
			Solve the equations by completing the square:
			\begin{enumerate}
				\item \( x^2 + 6x + 2 = 0 \)
				\item \( x^2 - 4x - 5 = 0 \)
				\item \( x^2 + 8x + 3 = 0 \)
				\item \( x^2 - 10x + 4 = 0 \)
			\end{enumerate}
		\end{alertblock}
	\end{frame}
	
	\begin{frame}
		\frametitle{Solutions – Completing the Square}
		\setbeamercolor{block body}{bg=green!20, fg=black}
		\begin{block}{Solutions}
			\begin{enumerate}
				\item \( x^2 + 6x + 2 = 0 \) \\
				{\color{blue} \( (x+3)^2 - 7 = 0 \)} \quad \( \Rightarrow x = -3 \pm \sqrt{7} \)
				
				\item \( x^2 - 4x - 5 = 0 \) \\
				{\color{blue} \( (x-2)^2 - 9 = 0 \)} \quad \( \Rightarrow x = 2 \pm 3 \Rightarrow x = 5, -1 \)
				
				\item \( x^2 + 8x + 3 = 0 \) \\
				{\color{blue} \( (x+4)^2 - 13 = 0 \)} \quad \( \Rightarrow x = -4 \pm \sqrt{13} \)
				
				\item \( x^2 - 10x + 4 = 0 \) \\
				{\color{blue} \( (x-5)^2 - 21 = 0 \)} \quad \( \Rightarrow x = 5 \pm \sqrt{21} \)
			\end{enumerate}
		\end{block}
	\end{frame}
	
	\begin{frame}
		\frametitle{From Completing the Square to the General Formula}
		\begin{block}{Introduction}
			We’ve learned how to solve a quadratic equation by completing the square. Now let’s see how this method leads to a general formula that works for any quadratic equation of the form \( ax^2 + bx + c = 0 \).
		\end{block}
		
		\setbeamercolor{block title}{fg=black, bg=blue!20}
		\setbeamercolor{block body}{fg=black, bg=blue!10}
		\begin{block}{General Formula}
			By completing the square, we derive the general formula to solve any quadratic equation:
			\[
			x = \frac{-b \pm \sqrt{b^2 - 4ac}}{2a}
			\]
			This is the \textbf{quadratic formula}, which gives the solutions of \( ax^2 + bx + c = 0 \).
		\end{block}
	\end{frame}
	
	\begin{frame}
		\frametitle{The General Formula}
		\begin{block}{}
			Consider the equation \( ax^2 + bx + c = 0 \), with \( a \neq 0 \). Divide by \( a \):
			\[
			x^2 + \frac{b}{a}x + \frac{c}{a} = 0
			\]
			Complete the square:
			\[
			x^2 + \frac{b}{a}x + \left(\frac{b}{2a}\right)^2 - \left(\frac{b}{2a}\right)^2 + \frac{c}{a} = 0
			\]
			\[
			\Rightarrow \left(x + \frac{b}{2a}\right)^2 = \frac{\textcolor{blue}{b^2 - 4ac}}{4a^2}
			\]
			Before taking the square root, we analyze the right side. Since \( 4a^2 > 0 \), we focus on the term \( \textcolor{blue}{b^2 - 4ac} \), which is called the \textbf{discriminant} and denoted by \( \Delta \).
		\end{block}
	\end{frame}
	
	\begin{frame}
		\frametitle{The General Formula: The Discriminant and Solutions}
		\begin{block}{}
			Now let’s examine the discriminant \( \Delta = b^2 - 4ac \) and the resulting cases:
			
			\begin{itemize}
				\item \textbf{If} \( \Delta > 0 \): The equation has two distinct real solutions.
				\[
				x_{1,2} = \frac{-b \pm \sqrt{\Delta}}{2a}
				\]
				
				\item \textbf{If} \( \Delta = 0 \): The equation has one real (repeated) solution.
				\[
				x = -\frac{b}{2a}
				\]
				
				\item \textbf{If} \( \Delta < 0 \): The equation has no real solutions.
			\end{itemize}
		\end{block}
	\end{frame}
	
	\begin{frame}
		\frametitle{Exercises}
		\begin{alertblock}{Solve the following equations:}
			\begin{enumerate}
				\item \( 2x^2 + 3x = -1 \)
				\item \( -x^2 = x - 12 \)
				\item \( 2x^2 - 5 = -4x \)
				\item \( 2(x - 1) = 3x^2 \)
				\item \( x(x - 2) = 15 \)
				\item \( x(x - 5) = 24 \)
				\item \( 2(x + 3)(x + 1) = x + 15 \)
			\end{enumerate}
		\end{alertblock}
	\end{frame}
	
	\begin{frame}
		\frametitle{Solutions}
		\begin{exampleblock}{Solutions}
			\begin{enumerate}
				\item \( \{-\frac{1}{2}, -1\} \)
				\item \( \{-4, 3\} \)
				\item \( \left\{-1 \pm \frac{\sqrt{14}}{2}\right\} \)
				\item No real solutions
				\item \( \{-3, 5\} \)
				\item \( \{-3, 8\} \)
				\item \( \left\{-\frac{9}{2}, 1\right\} \)
			\end{enumerate}
		\end{exampleblock}
	\end{frame}
	
	\begin{frame}
		\begin{alertblock}{Advanced Equation}
			\[
			x^2 - (\sqrt{3} - \sqrt{2})x - \sqrt{6} = 0
			\]
		\end{alertblock}
	\end{frame}
	
	\begin{frame}
		\frametitle{Solving the Equation}
		\fontsize{9}{10}
		\begin{exampleblock}{Solution}
			Given:
			\[
			x^2 - (\sqrt{3} - \sqrt{2})x - \sqrt{6} = 0
			\]
			with \(a = 1\), \(b = -(\sqrt{3} - \sqrt{2})\), \(c = -\sqrt{6}\)
			
			Compute the discriminant:
			\[
			\Delta = b^2 - 4ac = (\sqrt{3} - \sqrt{2})^2 + 4\sqrt{6}
			\]
			\[
			b^2 = 3 + 2 - 2\sqrt{6}, \quad -4ac = 4\sqrt{6}
			\]
			\[
			\Delta = 5 - 2\sqrt{6} + 4\sqrt{6} = 5 + 2\sqrt{6} = \textcolor{blue}{(\sqrt{3} + \sqrt{2})^2}
			\]
			
			So:
			\[
			x = \frac{-b \pm \sqrt{\Delta}}{2} = \frac{\sqrt{3} - \sqrt{2} \pm \textcolor{blue}{(\sqrt{3} + \sqrt{2})}}{2}
			\]
			Thus:
			\[
			x_1 = \sqrt{3}, \quad x_2 = -\sqrt{2}
			\]
		\end{exampleblock}
	\end{frame}
	
	
	
\end{document}
